% Academic poster writing scaffold; not an event-issued template.
% Reviewed: 2026-10-01. Confirm physical dimensions and export rules with the event.
% Compile: pdflatex -no-shell-escape beamerposter_academic.tex
\documentclass[final]{beamer}
% Custom width/height values use centimetres. Do not resize after setting type.
\usepackage[size=a0,scale=1.24,orientation=portrait]{beamerposter}
\usepackage[utf8]{inputenc}
\usepackage{graphicx,booktabs}
\usepackage{qrcode}
\usetheme{default}
\definecolor{OIdarkblue}{RGB}{0,114,178}
\definecolor{OIvermillion}{RGB}{213,94,0}
\setbeamercolor{headline}{fg=white,bg=OIdarkblue}
\setbeamercolor{block title}{fg=white,bg=OIdarkblue}
\setbeamercolor{block body}{fg=black,bg=white}
\setbeamercolor{block alerted title}{fg=white,bg=OIvermillion}
\setbeamercolor{block alerted body}{fg=black,bg=white}
\setbeamerfont{title}{size=\VeryHuge,series=\bfseries}
\setbeamerfont{author}{size=\Large}
\setbeamerfont{institute}{size=\large}
\setbeamerfont{block title}{size=\Large,series=\bfseries}
\setbeamerfont{block body}{size=\normalsize}
\setbeamertemplate{navigation symbols}{}
\setbeamertemplate{footline}{}

\title{Your Research Title Here: A Concise and Descriptive Title}
\author{First Author\inst{1}, Second Author\inst{1,2}, Third Author\inst{2}}
\institute[shortinst]{%
\inst{1} Department of Science, University Name, City, State, Country\\
\inst{2} Institute of Research, Institution Name, City, Country}

\begin{document}
\begin{frame}[t]
% Explicit title area: merely declaring \title does not display it in a frame.
\begin{beamercolorbox}[wd=\textwidth,sep=1.4cm,center]{headline}
{\usebeamerfont{title}\inserttitle\par}
\vspace{1cm}
{\usebeamerfont{author}\insertauthor\par}
\vspace{0.6cm}
{\usebeamerfont{institute}\insertinstitute\par}
\end{beamercolorbox}
\vspace{1.5cm}

\begin{columns}[t,onlytextwidth]
\begin{column}{.48\textwidth}
\begin{block}{Research question}
[State the scientific question and why resolving it matters.]
\vspace{0.6cm}
\begin{itemize}
\item \textbf{Established evidence:} [Relevant verified background.]
\item \textbf{Gap:} [What remains unresolved.]
\item \textbf{Objective:} [The specific question tested here.]
\end{itemize}
\end{block}
\vspace{1cm}

\begin{block}{Methods}
\begin{itemize}
\item \textbf{Design and units:} [Design, sampling frame, and independent experimental units.]
\item \textbf{Samples:} [Actual sample sizes, exclusions, and relevant characteristics.]
\item \textbf{Measurements:} [Intervention, comparator, outcomes, and timing.]
\item \textbf{Analysis:} [Estimand, model, uncertainty, missing-data handling, and software versions actually used.]
\end{itemize}
\vspace{1cm}
\centering
\fbox{\parbox[c][12cm][c]{.88\linewidth}{\centering
Replace with a labelled study-design diagram.\\[0.8cm]
Show the actual sampling, measurement, and analysis sequence.}}
\end{block}
\vspace{1cm}

\begin{block}{Evidence and limitations}
[Explain controls, sensitivity checks, and the assumptions needed to interpret the result.]
\vspace{0.6cm}
\begin{itemize}
\item {[Principal source of uncertainty or bias.]}
\item {[Boundary on generalization or causal interpretation.]}
\item {[What the present evidence cannot establish.]}
\end{itemize}
\end{block}
\end{column}

\begin{column}{.48\textwidth}
\begin{block}{Main result}
[Write one conclusion supported by the displayed data.]
\vspace{1cm}
\centering
% Replace the framed placeholder with an actual figure, for example:
% \includegraphics[width=.95\linewidth]{figure1.pdf}
\fbox{\parbox[c][20cm][c]{.88\linewidth}{\centering
Insert the main data figure.\\[0.8cm]
Label quantities, units, group sizes, and uncertainty.\\[0.8cm]
Use redundant labels or shapes so color is not the only cue.}}\par
\vspace{0.7cm}
\raggedright
\noindent\textbf{Figure 1.} [Describe the actual comparison. Define the uncertainty bars,
independent sample size, and any statistical annotations. No results are supplied
by this scaffold.]
\end{block}
\vspace{1cm}

\begin{alertblock}{Take-home message}
[State what was learned, with the scope and uncertainty warranted by the study.]
\end{alertblock}
\vspace{1cm}

\begin{block}{References, acknowledgments, and access}
[Add essential verified references, funding, disclosures, and contributor credits.]\par
\vspace{1cm}
\textbf{Contact:} [email protected]\\
\textbf{Paper or materials:} [Short, verified public URL or access instructions.]
% Optional QR: replace example.invalid with the intended permitted destination,
% then enable and test the code on the final exported PDF and a printed crop.
% \qrcode[height=7cm]{https://example.invalid/replace-before-use}
\end{block}
\end{column}
\end{columns}
\end{frame}
\end{document}

% Inspect a full-size printed crop for text and figure readability.
% Confirm PDF page dimensions, embedded fonts, QR destination, and print profile.
% A compilation pass is not proof of accessibility or event compliance.
